% limbtrack: technical description and SOP text for KVT3-PROJ-01 (lower-extremity MVIC test day)
% Written 2026-10-26 for limbtrack 0.20.0, SOP v0.18 (2026-10-21) and the protocol draft of 2026-10-20.
% Compiles on its own (pdfLaTeX, Overleaf). To use it in the SOP: copy Part A as an appendix, replace the SOP steps with Part B,
% add Table C1 as the settings checklist, and apply Part D to the protocol. Open points are \tbd{...} (red), as in the SOP.
\documentclass[a4paper,10pt]{article}
\usepackage[utf8]{inputenc}
\usepackage[T1]{fontenc}
\usepackage{lmodern}
\usepackage[margin=2cm]{geometry}
\usepackage{booktabs,tabularx,longtable,array,enumitem,amsmath,xcolor}
\usepackage[hidelinks]{hyperref}

\newcommand{\lt}{\textsf{limbtrack}}
\newcommand{\ltversion}{0.20.0}
\newcommand{\code}[1]{\texttt{#1}}
\newcommand{\dg}{\ensuremath{^\circ}}
\newcommand{\tbd}[1]{\textcolor{red}{[TBD: #1]}}
\newcolumntype{L}{>{\raggedright\arraybackslash}X}
\newcolumntype{P}[1]{>{\raggedright\arraybackslash}p{#1}}
\setlist{nosep,leftmargin=*}

\title{\lt{} \ltversion: technical description and SOP text\\[2pt]\large KVT3-PROJ-01, lower-extremity MVIC test day}
\author{KVT3 project group (lower body)}
\date{2026-10-26 (draft, for SOP v0.19)}

\begin{document}
\maketitle

\noindent This document has four parts. \textbf{Part A} describes how \lt{} works, technically, for an appendix of the SOP and for
AnyBody Technology. \textbf{Part B} is replacement text for every SOP step that uses \lt{}. \textbf{Part C} is the list of settings to
check at the start of each session. \textbf{Part D} lists the sentences of the research protocol that disagree with \lt{} or with the SOP.
Everything refers to \lt{} \ltversion; the version is saved in every trial (\code{meta.json}).

\tableofcontents

%=====================================================================================================================================
\part*{Part A. How \lt{} works}
\addcontentsline{toc}{part}{Part A. How limbtrack works}

\section{Purpose and scope}
\lt{} is the group's camera-based tracking program. During each MVIC trial it (1) tracks felt-pad markers on bony landmarks in a live
side or front view, (2) computes the joint angles of Table~3 of the protocol in every picture and shows the deviation from the target
position to the operator, (3) records the load-cell force at 1000 samples per second on the same clock, (4) marks the start and end of the
contraction from the force and gives sound cues in the 1--5--1 tempo, and (5) saves the raw video, a video with the overlay, the angles,
the dot positions and the force of the trial. Its Review tab turns a trial into the numbers used in the analysis (angle during the hold,
peak force, resting force, moment arm, torque, checks) and exports one row per trial.

\lt{} measures in two dimensions, in the plane of the picture. It has \textbf{not yet been validated} against a goniometer or an
inclinometer on a person; the angle-meter readings of SOP step 9.3 are the first validation data (Section~\ref{sec:limits}).

\section{Hardware and software}
\begin{tabularx}{\linewidth}{@{}lL@{}}
\toprule
Part & What is used \\
\midrule
Computer & Windows 11 laptop, charged from the isolated socket (the DAQ's ground reference; without it the 50~Hz mains hum on the force was
  about 40 times larger). \\
Program & \lt{} \ltversion, Python 3.12 with PySide6 (window), OpenCV (pictures, video files), pyqtgraph (graphs), NI-DAQmx (force).
  Started with \code{Start.bat}. Every trial saves the version, a hash of the program code, the git commit, the library versions and a
  hash of \code{protocol.json}. \\
Camera & Phone camera on a tripod, connected by USB through DroidCam (not Wi-Fi). \lt{} opens the camera by name and asks for
  1280$\times$720 pixels at the camera's rate (about 30 pictures per second). Only one program can use the phone stream at a time. \\
Force & TAS501c S-type load cell (500~kg) $\rightarrow$ SIC-A2 amplifier $\rightarrow$ NI USB-6003, differential input \code{ai0},
  1000 samples per second. Only one program can hold the DAQ (close MATLAB and NI MAX). \\
Test definitions & \code{protocol.json}: the dots, the angles (Table~3), the tests with their positions and required neighbour angles
  (Table~2), the tolerances. Checked when the program starts. \\
\bottomrule
\end{tabularx}

\section{A session in the program}
The Live tab has a strip of ten steps in the order of a session: \emph{Participant, Camera, Privacy, Position, Dots, Ruler, Lock, Force,
Markers, Record}. A step is green when done and amber while it is still to do; \emph{Privacy, Ruler, Force} and \emph{Markers} are asked
every session (they stay amber until done or until ``not today'' is ticked). Under the strip a bar stays in view with the position (a
dropdown of the whole session plan with the number of saved trials), one button per dot and \emph{Redo all}. Keys: R record, Space lock,
Z zero, G reading taken now, S/E MVIC start/end by hand, D redo all dots, 1--9 redo one dot, N/B next/previous position, P pause.

\begin{itemize}
\item \textbf{Participant.} The participant code (P01, \dots; codes only) and the \textbf{tested leg}. The tested leg belongs to the
  code: a new code starts with no leg chosen, the step reads \emph{P05: leg?}, and Record asks \emph{Right leg / Left leg}. The leg is saved
  in every trial. A code kept from the previous session (within 12~h) must be confirmed before it counts.
\item \textbf{Session plan.} The 19 positions and the two belt-squat positions in the SOP's fixed order, with the trials saved per
  position for the participant. A position is ``done'' at two trials, or reads \code{3rd!} when the 10~\% check asks for a third.
\end{itemize}

\section{Markers and tracking}\label{sec:tracking}
\begin{itemize}
\item \textbf{Markers.} Black felt pads ($\geq$2~cm) on white sports tape, on the skin (\emph{Dot colour = Black dots}). The pelvis
  landmark of the hip angle is one pad on the iliac crest, midway between ASIS and PSIS, because ASIS and PSIS are seen edge-on by a
  side camera (\code{PELVIS\_LANDMARK.md}). Frontal tests add pads on both ASIS and the patella centre.
\item \textbf{Marking.} The operator clicks each dot once in the order asked; a click finds a dot up to about 12~pixels away (limit
  15~px). A dot can be marked again on its own (its number key removes its old mark at once) or all in turn (D removes every old mark and the lock at once).
\item \textbf{Tracking.} In every picture each dot is found again near its last position as a round, dark (or bright), right-sized
  blob; its centre is the dot position. All tracking constants are fixed and saved with every trial.
\item \textbf{Lost dots.} A dot that is not found is \emph{LOST}: it stays at its last good position, is searched for only within
  2$\times$ its diameter, and is taken back only on a round, right-sized, contrasting blob. No angle is computed from a lost dot and nothing
  is interpolated: the CSV has an empty cell. A dot seen only partly is flagged \emph{partial}.
\item \textbf{Label checks.} The dots must lie in the anatomical order (for example, the malleolus farther from the epicondyle than the
  fibular head; in frontal tests the tested ASIS on the side of the tested leg). A failing check shows \emph{CHECK LABELS}, blocks Lock, and is
  saved per picture (\code{label\_check}).
\item \textbf{Dots on one bone.} The distance between two dots on the same bone (fibular head--malleolus, 5th metatarsal base--head,
  trochanter--epicondyle, the two ASIS) is taken when they are marked. When it changes by more than 25~\%, a tracker has most likely jumped
  to something else: the angles using them show \emph{CHECK DOTS}, Lock is refused, and the frames are saved (\code{dots\_apart}). The angle
  value itself is still saved unchanged.
\end{itemize}

\section{Joint angles}\label{sec:angles}
Each angle is the angle between two ``arms'' in the picture, an arm being the vector from one dot to another, as the goniometer arms of
Table~3. The shown value is $\text{sign}\times\text{angle}+\text{offset}$ (formulas in Section~\ref{sec:angle-formulas}). Signed angles depend on which way the participant faces in the
picture (\emph{Participant faces}); the program warns when the dots show the other way (from the foot dots, or from a knee bent the wrong
way).

\begin{tabularx}{\linewidth}{@{}P{2.4cm}P{3.1cm}P{3.1cm}L@{}}
\toprule
Angle & Proximal arm & Distal arm & Zero and sign \\
\midrule
Hip & iliac crest pad $\rightarrow$ greater trochanter & trochanter $\rightarrow$ lateral epicondyle & 0 = thigh in line with the pelvis
  line; flexion positive, extension negative. Fulcrum: trochanter. \\
Knee & trochanter $\rightarrow$ lateral epicondyle & fibular head $\rightarrow$ lateral malleolus & 0 = straight; flexion positive,
  hyperextension negative. Fulcrum: epicondyle. \\
Ankle & malleolus $\rightarrow$ fibular head & 5th metatarsal base $\rightarrow$ head & 90\dg{} $-$ angle: 0 = foot square to the shank;
  dorsiflexion positive, plantarflexion negative ($-20$ = 20\dg{} plantarflexion). Fulcrum: malleolus. \\
Hip abduction (front view) & tested ASIS $\rightarrow$ other ASIS & tested ASIS $\rightarrow$ patella centre & angle $-$ 90\dg: 0 = thigh
  square to the line between the ASIS; abduction positive, adduction negative; the same for either leg. Fulcrum: tested ASIS. Used for
  both frontal tests: hip adduction is measured at 0 and 20\dg{} \emph{abduction}. \\
Hip abduction about the hip centre & as above & estimated hip joint centre $\rightarrow$ patella & shown and saved, never judged. The hip
  centre is computed from the two ASIS and the ruler scale (Harrington et al.\ 2007) \tbd{constants to be checked in the original paper}. \\
\bottomrule
\end{tabularx}

\medskip\noindent The saved angles are never smoothed. On screen they are shown as a 0.2~s moving average (display only).

\subsection{Formulas}\label{sec:angle-formulas}
\noindent This is the calculation exactly as \lt{} does it in every picture. Written out as a separate program from this text, it gives the
same value as \lt{} for all five angles of Table~3 on 2000 to 3000 random sets of dot positions each, for a participant facing either
way (no difference beyond rounding).

\paragraph{Dots and arms.} A dot is the centre of its marker, $(x,y)$ in picture pixels, with $x$ to the right and $y$ \emph{downwards}.
The arm from dot $A$ to dot $B$ is the vector
\begin{gather}
  \mathbf{v}_{AB}=\bigl(m\,(x_B-x_A),\ y_B-y_A\bigr),\label{eq:arm}\\
  m=\begin{cases}+1 & \text{the participant faces right in the picture,}\\ -1 & \text{the participant faces left.}\end{cases}\notag
\end{gather}
The factor $m$ mirrors the picture of a participant who faces left, so that a sign means the same whichever way the participant faces. It
makes no difference to the unsigned angles (ankle and front view).

\paragraph{The angle between two arms.} Let $\mathbf{v}_1\cdot\mathbf{v}_2=v_{1x}v_{2x}+v_{1y}v_{2y}$ and
$\mathbf{v}_1\times\mathbf{v}_2=v_{1x}v_{2y}-v_{1y}v_{2x}$, and let $\operatorname{atan2}(y,x)$ be the four-quadrant arctangent. Then
\begin{gather}
  \theta=\operatorname{atan2}\bigl(\lvert\mathbf{v}_1\times\mathbf{v}_2\rvert,\ \mathbf{v}_1\cdot\mathbf{v}_2\bigr)\in[0\dg,180\dg],\label{eq:theta}\\
  \theta_s=\operatorname{atan2}\bigl(\mathbf{v}_1\times\mathbf{v}_2,\ \mathbf{v}_1\cdot\mathbf{v}_2\bigr)\in[-180\dg,180\dg].\label{eq:thetas}
\end{gather}
$\theta$ is the ordinary angle between the two arms (0\dg{} when they point the same way, 180\dg{} when they point opposite ways). $\theta_s$ is
the same angle with a sign: positive when the turn from $\mathbf{v}_1$ to $\mathbf{v}_2$ is clockwise on the screen. If an arm has length zero
the angle does not exist and the cell in the CSV is empty.

\paragraph{The shown value.}
\begin{equation}\label{eq:alpha}
  \alpha=s\,\theta+c\quad(\text{unsigned angles}),\qquad\qquad \alpha=s\,\theta_s+c\quad(\text{signed angles}),
\end{equation}
with the sign $s$ and the offset $c$ of the table below. They are set in \code{protocol.json} and saved with every trial. The letters
are: $C$ iliac crest pad, $G$ greater trochanter, $E$ lateral femoral epicondyle, $F$ fibular head, $M$ lateral malleolus, $B$ 5th
metatarsal base, $D$ 5th metatarsal head, $A$ tested ASIS, $O$ other ASIS, $P$ patella centre, $J$ estimated hip joint centre.

\medskip
\begin{tabularx}{\linewidth}{@{}P{2.5cm}cclccL@{}}
\toprule
Angle & $\mathbf{v}_1$ & $\mathbf{v}_2$ & Type & $s$ & $c$ & Shown value $\alpha$ \\
\midrule
Hip & $\mathbf{v}_{CG}$ & $\mathbf{v}_{GE}$ & signed & $-1$ & $0$ & $-\theta_s$ (flexion positive) \\
Knee & $\mathbf{v}_{GE}$ & $\mathbf{v}_{FM}$ & signed & $+1$ & $0$ & $\theta_s$ (flexion positive) \\
Ankle & $\mathbf{v}_{MF}$ & $\mathbf{v}_{BD}$ & unsigned & $-1$ & $90\dg$ & $90\dg-\theta$ (dorsiflexion positive) \\
Hip abduction, front view & $\mathbf{v}_{AO}$ & $\mathbf{v}_{AP}$ & unsigned & $+1$ & $-90\dg$ & $\theta-90\dg$ (abduction positive) \\
Hip abduction about the hip centre & $\mathbf{v}_{AO}$ & $\mathbf{v}_{JP}$ & unsigned & $+1$ & $-90\dg$ & $\theta-90\dg$ (shown and saved, never judged) \\
\bottomrule
\end{tabularx}

\medskip\noindent Why the signs are as they are:
\begin{itemize}
\item \textbf{Hip and knee.} For a participant facing right, both arms point down the leg when the leg is straight ($\theta_s=0$). Bending the hip
  swings the thigh forward, which is anticlockwise on the screen ($\theta_s<0$, hence $s=-1$); bending the knee swings the shank
  backwards, which is clockwise ($\theta_s>0$, hence $s=+1$). In both, flexion is positive. If the \emph{Participant faces} setting is
  wrong, a flexion reads as the same number with a minus sign; the facing warning above catches this, and Lock is refused until the
  setting is corrected.
\item \textbf{Ankle.} $\mathbf{v}_{MF}$ runs up the shank and $\mathbf{v}_{BD}$ along the foot towards the toes, so a foot square to the
  shank gives $\theta=90\dg$. Toes towards the shin make $\theta$ smaller (so $\alpha>0$, dorsiflexion); toes pointing away make it larger
  ($\alpha<0$, plantarflexion; $-20$ is 20\dg{} plantarflexion).
\item \textbf{Hip abduction (front view).} $\mathbf{v}_{AO}$ runs along the line between the two ASIS, so a thigh hanging square to that line
  gives $\theta=90\dg$. The knee moving away from the other leg makes $\theta$ larger ($\alpha>0$, abduction), towards it smaller
  ($\alpha<0$, adduction). Being unsigned, it is the same for either leg and for a mirrored picture.
\end{itemize}

\paragraph{Estimated hip joint centre (front view).} The point $J$ is calculated from the two ASIS pads, not clicked. Let $T$ and $O$ be the
centres of the tested and the other ASIS pad (pixels), $k$ the scale of the ruler in pixels per millimetre, $w=\lVert T-O\rVert$ and
$PW=w/k$ the distance between the two ASIS in millimetres. Then
\begin{gather}
  J=\tfrac12\,(T+O)+k\,\bigl(d_{\text{side}}\,\mathbf{u}+d_{\text{down}}\,\mathbf{n}\bigr),\label{eq:hipcentre}\\
  d_{\text{side}}=0.33\,PW+7.3\ \text{mm},\qquad d_{\text{down}}=0.30\,PW+10.9\ \text{mm},\notag
\end{gather}
where $\mathbf{u}=(T-O)/w$ points from the other ASIS to the tested ASIS and $\mathbf{n}$ is the unit vector square to $\mathbf{u}$ with
$\mathbf{n}\cdot(P-T)>0$, that is, pointing to the side of the line $TO$ where the patella is (``down'' in the pelvis' own frame, so a tilted
camera or a participant lying on one side does not matter). The constants are those of Harrington et al.\ (2007), equation set~1
\tbd{constants to be checked in the original paper}. Without a ruler scale $J$ is not calculated and the hip-centre angle is empty.

\paragraph{Worked example} (participant faces right; picture pixels): $G=(300,200)$, $E=(500,200)$, $F=(500,220)$, $M=(500,420)$,
$B=(480,440)$, $D=(580,450)$, $C=(300,40)$.
\begin{itemize}
\item \textbf{Knee:} $\mathbf{v}_{GE}=(200,0)$, $\mathbf{v}_{FM}=(0,200)$, cross product $40000$, dot product $0$, so
  $\theta_s=\operatorname{atan2}(40000,0)=90.00\dg$ and $\alpha_{\text{knee}}=90.00\dg$ (the shank is vertical, the thigh horizontal).
\item \textbf{Hip:} $\mathbf{v}_{CG}=(0,160)$, $\mathbf{v}_{GE}=(200,0)$, cross product $-32000$, dot product $0$, so $\theta_s=-90.00\dg$ and
  $\alpha_{\text{hip}}=+90.00\dg$.
\item \textbf{Ankle:} $\mathbf{v}_{MF}=(0,-200)$, $\mathbf{v}_{BD}=(100,10)$, cross product $20000$, dot product $-2000$, so
  $\theta=\operatorname{atan2}(20000,-2000)=95.71\dg$ and $\alpha_{\text{ankle}}=90\dg-95.71\dg=-5.71\dg$ (5.7\dg{} plantarflexion: the toes
  point slightly down).
\end{itemize}
\noindent The same numbers come out of \lt{} for these dot positions (knee $90.00$, hip $90.00$, ankle $-5.71$). Picture pixels are used as
they are: lens distortion and perspective are not corrected, so each angle is that of the limb projected on the picture plane and is
right only for a limb moving in a plane square to the camera (Section~\ref{sec:limits}).

\section{Judging the position, and Lock}
\begin{itemize}
\item \textbf{Tested joint:} judged against the target position $\pm$ the tolerance (5\dg{} in the study; green inside, red outside).
\item \textbf{Adjacent joints with a value in Table~2} (hip 90\dg{} in the knee tests, knee 0\dg{} or 90\dg{} in the hip and ankle tests):
  judged against that value $\pm$ 5\dg{} from the start.
\item \textbf{Lock} stores the other angles (0.2~s average, unsmoothed) as the reference, so drift during the trial shows. Lock is refused
  when a dot is missing or lost, a label check fails, dots on one bone moved apart, the picture faces the other way, or (frontal tests) no
  leg is chosen. Lock does \emph{not} require the tested joint to be within tolerance (a set-up a few degrees beyond the target, SOP 9.11,
  can be locked).
\item \textbf{Record} asks before recording when the position is not locked, the main angle cannot be measured, the force has not been
  zeroed since the start, or no leg is chosen.
\end{itemize}

\section{Scale, strap and moment arm}
\begin{itemize}
\item \textbf{Ruler (scale).} Two fixed pads 100~cm apart, clicked once, give pixels per centimetre. The scale is valid only at the
  pads' distance from the camera (the pads are at the depth of the joint; for the frontal tests at the depth of the ASIS). The pads are not
  tracked: mark them again whenever the camera or the bench moves. With the scale, dots under 2~cm are flagged.
\item \textbf{Strap and anchor dots} (sagittal tests). The line of pull is the line from the strap dot to the anchor dot. \lt{} saves the
  pull angle (the angle between the segment and the line of pull, and its difference from 90\dg) and the \textbf{moment arm}: the
  perpendicular distance from the fulcrum dot of the tested joint to the line of pull, in pixels and in centimetres. Surface landmarks
  are not exact joint centres (protocol 2.5).
\item \textbf{Frontal tests:} no strap or anchor dots and no moment arm; the ASIS-to-strap distance is measured with a tape (SOP 9.9)
  \tbd{pivot for the frontal moment arm: ASIS or estimated hip centre, with AnyBody}. The ASIS-to-ASIS tape measure (mm) is compared with
  the scale; a difference above 8~\% is warned.
\end{itemize}

\section{Force}\label{sec:force}
\subsection{Chain, calibration and zero}
The DAQ reads the amplifier output in volts at 1000 samples per second; every sample is saved as measured (\code{force.csv}). Newtons are
computed as $F = (V - V_0)\times k$ with $k = 526.95$~N/V, the regression of 2026-10-13 (0--250~kg up and down, 24 runs, $R^2 = 0.999996$),
and $V_0$ the zero. \textbf{Zero} (Z) is taken with the strap slack and the chain hanging in the test direction: it removes the
amplifier offset and what the hardware itself puts on the load cell. A zero is accepted when its 1~s reading does not drift (halves within
10~mV) and its noise is at most 30~mV; it is saved with each trial and forgotten when the program closes. Review warns for forces beyond
the heaviest calibration load (250~kg).

\subsection{Peak, smoothing and timing}
Force and pictures share one clock (the DAQ clock is aligned to it). The live and saved peak (\code{peak\_n}) is the highest 20~ms moving
average, one 50~Hz mains cycle, so hum is not counted as force; the highest single sample is kept as \code{peak\_raw\_n}. Review uses a
100~ms moving average by default (a multiple of 20~ms removes the hum exactly). The delay between a picture being taken and being stamped
(\emph{picture delay}, typically 0.1--0.3~s) is not yet measured \tbd{tap test}; Review can correct for it.

\subsection{Resting force}\label{sec:rest}
The \textbf{push} begins where the smoothed force was last within 10~\% of its rise above the low level before it (the 10th percentile
of the force before it passed half of its peak). The \textbf{resting force} is the \emph{median of the force from 3.0 to 0.5~s before the
push}. With less than 0.5~s of force before the push there is no resting force; with a spread (SD) above 10~\% of the rise (the strap was
loaded or moving), \lt{} says so after the trial and in Review. Since 0.24.0 a short rest is no longer reported: the resting force is taken
\emph{before} the trial, with the participant relaxed in the strap in the position, by pressing Zero~2 (X); \lt{} reminds about it after every
change of position (a reminder, not a stop; Section~\ref{sec:reminders}) \tbd{group: Zero 2 replaces the SOP's 3~s of rest after Record}.
The angle at rest is taken from the same stretch. Export columns: \code{force\_rest}, \code{force\_rest\_s} (seconds
of rest), \code{force\_peak\_above\_rest}. (Until 0.19.0 the resting force was the second before the window \emph{force above half of
peak}, which holds about half of the ramp: on the trials of 2026-10-15 the peak above rest came out 6--16~\% too low.)

\subsection{Automatic MVIC markers}
With \emph{Mark the MVIC automatically from the force} on, Record arms a detector: the force is averaged over 20~ms; its rest level is the
median from 3.0 to 0.5~s before the present moment; the start is called when the force differs from that level by at least the threshold
(25~N on the lab laptop \tbd{confirm the value}; raised to 6 times the rest noise when the rest is unsteady), has moved by at least 40~\% of
the threshold within 0.25~s, and both hold for 70~ms. The start marker is written about 0.1--0.2~s after the force began to rise (the
estimate of the true onset is saved), the end marker 7~s later (1--5--1), and the recording stops by itself 2~s after the end marker. S and
E by hand take over. The markers only mark time; they never change the force or an angle.

\subsection{Sound cues}
Record: one beep. Ramp registered (the start marker): two rising beeps. One second later: one long high beep (hold). During the hold: a
short tick every second. One second before the end: two falling beeps (relax). End marker: one low beep. The times follow the force onset,
so the 5~s hold is timed by the program. The cue settings are saved with each trial because participants hear them.

\section{Checks after each trial}
When a trial is saved, the record line and a message that does not block (R starts the next trial) name every problem; each is also
written to the session log (\code{session\_log.csv}):
\begin{itemize}
\item \textbf{No force rise:} the force moved less than 60~N between the MVIC markers (strap not on the load cell, amplifier off, cable out).
\item \textbf{CHECK DOTS} / \textbf{CHECK LABELS} in some pictures (a tracker jumped, or two dots swapped labels).
\item \textbf{Rest:} the force was not steady before the push (Section~\ref{sec:rest}). A short rest is not reported any more.
\item \textbf{10~\% check (SOP 10.6):} all saved trials of the same participant, test and target are compared. The total peak force is
  used (the export's \code{force\_peak}); the spread is $(\max-\min)/\text{mean}$. Two trials more than 10~\% apart: \emph{THIRD TRIAL
  NEEDED} (and \code{3rd!} in the plan). Three trials: judged on the closest two; none within 10~\%: note it in the deviation log. The
  peak above rest is shown beside it for information.
\end{itemize}

\section{Reminders instead of stops, and the dots of each test}\label{sec:reminders}
Since 0.24.0 the steps run Camera, Force, Markers, Participant, Position, Privacy, Ruler, Dots, Lock, Record. After every change of
position (and at the first look, for privacy) rows appear under the instruction line, each with the buttons that answer it:
\begin{itemize}
\item \textbf{Privacy:} \emph{No face} (key V), \emph{Mask}, \emph{Crop}; an amber banner across the live picture (never in a saved video).
\item \textbf{Ruler:} \emph{Fine}, or \emph{Redo} (key C, also the \emph{Ruler (C)} button beside the position).
\item \textbf{Resting force:} \emph{Zero 2} (key X), or \emph{Not here}.
\end{itemize}
They do not stop Record. What was still open when Record was pressed is written to the session log and to \code{meta.json}
(\code{reminders\_open\_at\_start}); a trial recorded with the privacy not looked at again, the position not locked, a doubtful main angle or
the force not zeroed says so with its warnings after it is saved. Still stopping a recording: no video, paused, a participant code that is not
P01\dots, notes that could identify someone, a crop that does not fit the picture; and the quick questions \emph{same participant?} and
\emph{which leg?}.

\textbf{Dots of each test.} The foot (the 5th metatarsal base and head, and the ankle angle) is tracked only in the ankle tests and the belt
squat \tbd{group: the foot and belt-squat analysis may move to post}; the knee and hip tests use the five dots iliac crest, trochanter,
epicondyle, fibular head and malleolus. One word per test in \code{protocol.json} brings the foot back.

\section{Recording, files and privacy}
Each trial is saved in \code{recordings/<code>/} as \code{<code>\_<NN>\_<test>\_<target>deg\_T<n>\_<date>-<time>} plus the file
type, for example \code{P01\_03\_knee\_extension\_110deg\_T2\_20261026-101712\_raw.mp4}: NN is the position number, the No.\ of the SOP's test
order (01--21, as shown in the session plan; 00 for a typed target that is not a plan position), and n the trial number at that position.
A participant's files therefore sort in the order of the session. Both numbers are also saved in \code{meta.json} and exported
(\code{position\_number}, \code{trial\_number}).

\medskip

\begin{tabularx}{\linewidth}{@{}lL@{}}
\toprule
File & Content \\
\midrule
\code{\_raw.mp4} & the camera picture (cropped and masked for privacy) \\
\code{\_overlay.mp4} & the same with dots, angles and the header drawn (optional) \\
\code{\_data.csv} & one row per picture: time, every dot's position (empty while lost), quality and partial flag, every angle with its
  reference, deviation and within-tolerance flag, \code{all\_ok}, markers, scale, mean force since the previous picture (\code{force\_n},
  \code{force\_v}), pull angle and moment arm, \code{label\_check}, \code{dots\_apart}. Angles are unsmoothed. \\
\code{\_force.csv} & every force sample (time, volts, newtons) \\
\code{\_meta.json} & everything needed to reproduce the trial: participant, tested leg, test, target, tolerances, the angle definitions,
  the settings at the start, camera, privacy, zero and calibration, automatic markers and their times, cue settings, checks, frame rate and
  timing, reference readings, program version and hashes. Written when the recording starts and replaced at the end. \\
\bottomrule
\end{tabularx}

\medskip\noindent Rows are written as they come and the videos are encoded on their own threads; pictures are never dropped. The crop
and black masks apply to both saved videos only; tracking and the CSV use the whole picture (in full-picture pixels). Recording does not
start when a tracked dot is within 30~px of the crop edge or inside a mask. The Notes field is saved in every \code{meta.json} and must never
hold a name.

\section{Review and the numbers for the analysis}
\begin{itemize}
\item \textbf{Window.} The export uses the stretch around the peak where the smoothed force is above half of its peak (\emph{force above
  half of peak}); without force, the MVIC markers.
\item \textbf{Plateau angle} (\code{plateau\_angle\_deg}): the mean of the tested joint's angle over the middle 60~\% of the window.
  \textbf{Angle at peak force} (\code{angle\_at\_peak\_force\_deg}): the mean angle of the pictures within 0.1~s of the peak (after the
  picture delay). \tbd{which of the two goes to AnyBody (SOP 10.7)}
\item \textbf{Angle rest to hold} (\code{angle\_rest\_to\_hold\_deg}): plateau angle minus the angle at rest (the joint's movement under
  load: slack, padding, stretch).
\item \textbf{Force:} peak, plateau mean, resting force, peak above rest, onset, time from 10 to 90~\% of the rise, steadiness of the
  hold, hum as a share of the rise.
\item \textbf{Torque:} smoothed force $\times$ moment arm, when the force is calibrated and the ruler, strap and anchor dots are present.
\item \textbf{Reference readings:} the angle-meter value typed for each joint, the program's angle at the moment G was pressed (0.2~s
  average), their difference, and over all trials the mean difference and 95~\% limits of agreement (Bland and Altman).
\item \textbf{Export all trials of this participant} writes one row per trial (\code{session\_summary.csv}) with the numbers above, the
  checks and the 10~\% columns.
\end{itemize}

\section{Using the force in AnyBody: zero versus resting force}\label{sec:anybody}
Moments about the tested joint, positive in the direction the participant pushes. With the strap force $F$ at moment arm $d$, static
equilibrium gives, at rest (muscles relaxed) and at the peak:
\begin{align*}
F_\text{rest}\,d_\text{rest} &= M_\text{gravity}(\theta_\text{rest}) + M_\text{passive}(\theta_\text{rest}),\\
F_\text{peak}\,d_\text{peak} &= M_\text{active} + M_\text{gravity}(\theta_\text{hold}) + M_\text{passive}(\theta_\text{hold}),
\end{align*}
where $M_\text{gravity}$ is the moment of the weight of the limb beyond the joint (with the strap and any hardware hanging on it) and
$M_\text{passive}$ that of stretched muscles, tendons and capsule. The zero removes only the amplifier offset and the load-cell hardware;
it does not remove either of these.

\begin{enumerate}
\item \textbf{The exact set-up modelled in AnyBody, with gravity: use the total zeroed force} ($F_\text{peak}$, column \code{force\_peak}).
  The model adds $M_\text{gravity}$ itself. \textbf{Subtracting the resting force as well would remove gravity twice.} The simple muscle
  model has no passive elements, so its scaled strength then includes $M_\text{passive}$: small in mid-range positions, larger where muscles
  are stretched (for example knee flexion 20\dg{} with the hip at 90\dg, hip extension at 45\dg{} with the knee straight).
\item \textbf{Peak above rest} ($F_\text{peak}-F_\text{rest}$, column \code{force\_peak\_above\_rest}) estimates $M_\text{active}$ alone.
  It is valid only when the model's MVIC load case does \emph{not} include the limb's gravity, and only when the posture is the same at rest
  and at the peak ($\theta_\text{rest}\approx\theta_\text{hold}$, $d_\text{rest}\approx d_\text{peak}$, same direction of pull).
\item \textbf{To exclude passive tension while modelling gravity:} use $F_\text{peak} - (F_\text{rest} - F_\text{rest}^\text{model})$,
  where $F_\text{rest}^\text{model}$ is the strap force the AnyBody model needs to hold the relaxed limb in the rest posture (gravity only,
  muscles off). The difference $F_\text{rest} - F_\text{rest}^\text{model}$ is the passive tension and any pre-tension of the set-up.
\end{enumerate}

\noindent\textbf{Rule: subtract the resting force or model the limb's gravity, never both.} Whatever is chosen, the resting force is
usable only when the participant was relaxed (\lt{} flags a rest that was not steady), at rest for at least 3~s, and in the same posture as
during the hold (\code{angle\_rest\_to\_hold\_deg}: the knee moved about 6\dg{} under load in the first force trials, and up to 19\dg{} is
reported \cite{tsaopoulos}). Recommendation: give AnyBody the total zeroed force with the hold angle and the pull direction at the peak,
and use the resting force as a check: $F_\text{rest}\,d$ should be close to the model's gravity moment in the rest posture; a large
difference means passive tension, pre-tension or a participant who was not relaxed \tbd{decide with AnyBody whether passive tension is
removed (option 3)}. Hardware between the load cell and the limb (strap, carabiner, in vertical pulls possibly the cell) loads the limb but
is not in the reading: about 10--15~N at most; include it in the model as a small force at the strap or state it.

\section{Settings remembered per computer}
\lt{} remembers its set-up in \code{settings.json} on each computer (tolerances, dot colour, markers, cues, ruler length, camera). The
participant code and the tested legs are remembered for 12~h only. The zero, the marked dots, the lock and the ``not today'' ticks are
never remembered. Because a remembered value can differ from the SOP's, the operator checks the values in Table~C1 at the start of each
session.

\section{Known limitations and validation status}\label{sec:limits}
\begin{itemize}
\item Not yet validated against a goniometer or inclinometer on a person; the readings of SOP 9.3 are the validation data (target:
  mean difference within 0.5\dg, limits of agreement within $\pm$1\dg, which a hand-held instrument read to 1\dg{} cannot demonstrate).
\item 2D: a segment turned out of the picture's plane is shortened and its angle changes; the camera must be square to the plane of
  movement and level with the joint.
\item Skin and pads move over the bone; the trochanter in particular in hip flexion. Seated at a nominal hip 90\dg{} the hip angle read
  37--68\dg{} in earlier sessions \tbd{hip-first angle-meter readings in the pilot decide whether the pelvis line or the positioning is off}.
\item In the ankle tests the foot arm (5th metatarsal base to head) is short in the picture (33--69~px on 2026-10-22): one pixel of
  error is about 1\dg. Frame the ankle closely and keep both foot dots on skin, clear of the strap.
\item The estimated hip centre uses constants taken from a secondary source; the hip-centre angle is never judged.
\item The picture delay is not measured; the frame rate is saved with every trial and should be close to 30 pictures per second.
\end{itemize}

%=====================================================================================================================================
\part*{Part B. SOP steps that use \lt{} (replacement text)}
\addcontentsline{toc}{part}{Part B. SOP steps that use limbtrack}

\noindent Step numbers are those of SOP v0.18. Text in this part replaces the step; steps not listed stay as they are.

\begin{longtable}{@{}P{2.2cm}p{14.3cm}@{}}
\toprule
Step & Text \\
\midrule
\endhead
1.3 & Laptop with \lt{} \ltversion{} \tbd{the version frozen after the pilot}; the same version is used for every participant and is not
  updated during data collection (it is saved in every \code{meta.json}; write it on the session sheet). Phone camera on a tripod,
  connected to the laptop by USB with DroidCam, not over Wi-Fi. Only one program can use the phone stream at a time. \\
3.4 & Start \lt{}, choose NI DAQ, press Connect, and check that the force read-out responds to a light pull. Choose the study's
  recordings folder (Record section) \tbd{name of the folder}: a new, empty folder, never the folder with the pilot trials. \\
3.5 & Set up the camera according to the Camera set-up list (stabilisation off, zoom 1x, focus, exposure and white balance locked in
  the phone app) and run Check the camera (5~s). Check the settings in Table~C1. \\
4.5 & Determine the dominant leg by asking: ``Which leg do you kick a ball with?'' The dominant leg is the tested leg. In \lt{}, type the
  participant code and choose the tested leg in the Participant section; it is kept for this code and saved in every trial (Record asks
  for it if it is not chosen). \\
9.2 & Set the adjacent joint to its required angle (Table~2). \lt{} judges the adjacent joint from the start. \tbd{until the pilot has
  decided the hip 90\dg{} question, the angle meter decides the adjacent hip and knee angles; \lt{}'s colour for them is not used} \\
9.3 & Cross-check the joint angle at rest with the electronic angle meter, with its arms aligned as in Table~3. In \lt{}, open
  Reference reading and choose Goniometer (the A-Digit 50 has two arms) and resolution 0.1. Keep the participant still, press G once at the
  first reading, take up to five readings and type them in the joint's box separated by a space (they are averaged against that moment).
  Pressing G again replaces the moment. A reading belongs to this participant and position only. In the pilot, a reading is taken at every
  position, the hip first at the seven positions with the hip at 90\dg; in data collection, at the first position of each test. \\
9.5 & (unchanged) The scale pads must be at about the joint's distance from the camera; type 100~cm. Mark them again if the camera or the
  bench is moved. Frontal tests: pads on a board in the frontal plane at the depth of the ASIS. \\
9.6 & In \lt{}: choose the position in the dropdown under Lock/Record (N = next position), check \emph{Participant faces} (side view: it
  changes between seated, prone and supine) or the tested leg (front view), choose Black dots, mark the dots, position the joint (green
  within tolerance), and lock (Space). Check on screen that the tolerance for the tested joint reads 5.0 and for the other joints 5.0. A
  deviation within the tolerance is accepted, because the measured angle during the hold is recorded and used in AMS. \\
9.7 & In the sagittal tests, put a dot on the flat side of the strap facing the camera (not on the buckle) and one on the anchor, at least
  5~cm from any landmark dot. Landmark dots stay on the skin and clear of the strap edge (hip tests: the strap just proximal to the
  condyles is next to the epicondyle dot; plantarflexion: the strap at the metatarsal heads is at the 5th metatarsal head dot). \\
9.9 & \lt{} tracks the fulcrum dot and the strap dot and gives the moment arm (the perpendicular distance from the fulcrum to the line
  from the strap dot to the anchor dot) in cm. Check that both dots are tracked before recording. Frontal tests: no moment arm; measure the
  ASIS-to-strap distance with a tape and write it on the session sheet. \\
9.10 & Zero the load cell (Z) with the strap slack and the chain hanging in the test direction; zero again after any restart of \lt{}.
  Then attach the strap. With the participant in position and relaxed in the strap, press Zero~2 (X): this is the resting force of the
  position (\lt{} reminds about it after every change of position; it does not stop Record) \tbd{group: replaces the 3~s of rest after
  Record}. \lt{} also finds the rest in the recording (the median force from 3.0 to 0.5~s before the push) and says after the trial when
  the force was not steady. Both the total force and the force minus the resting force are exported; which one AnyBody gets
  follows Part~A, Section~\ref{sec:anybody}. \\
9.11 & (unchanged; \lt{} does not require the tested joint to be within tolerance to lock, so a set-up a few degrees beyond the target
  can be locked). \\
10.1 & DR checks that \emph{Mark the MVIC automatically from the force} is on and presses Record while the participant is at rest. \\
10.2--10.3 & TL gives the verbal cue with \lt{}'s beeps: ``Ready'' after the record beep \tbd{the 3~s of rest after Record is replaced by the resting force, Zero 2}, ``build up'' (two rising beeps
  follow when the ramp is seen), ``max, push, push, push'' from the long beep (a tick every second), ``and relax'' on the two falling beeps.
  TL watches for compensation (trunk lean, pelvic tilt, lifting off the support). \\
10.4 & The recording stops by itself 2~s after the end marker. DR reads the record line and the message: \emph{no force rise} $\rightarrow$
  check chain, load cell and amplifier and repeat; \emph{CHECK DOTS} or \emph{CHECK LABELS} $\rightarrow$ mark the dot again (its number
  key, or D for all) and repeat; \emph{force not steady before the push} $\rightarrow$ note it (repeat if the resting force is needed). A trial with lost
  dots or a clear compensation is repeated. If the hold is outside $\pm$5\dg, this is noted, and the set-up is adjusted before the next
  trial. \\
10.6 & After the second trial \lt{} gives the 10~\% check of the position (total peak force, spread over the mean of the two). \emph{THIRD
  TRIAL NEEDED} $\rightarrow$ rest 60~s and record a third trial. After a third trial it judges the closest two; if no two are within
  10~\%, note it in the deviation log. No more than three trials are performed. \\
10.7 & Record on the session sheet: trial numbers, peak force, the measured joint angle during the hold \tbd{plateau mean or angle at
  peak force}, and any remarks. \\
11.6 & In \lt{}, choose the belt squat position (knee flexion 90\dg, then 60\dg; 0\dg{} = straight knee), mark the dots and lock. The
  knee and hip angles are recorded by \lt{}. \\
13.1 & Check that each trial in the study's recordings folder \code{<code>/} has \code{\_raw.mp4}, \code{\_overlay.mp4},
  \code{\_data.csv}, \code{\_meta.json} and \code{\_force.csv}. The file names begin with the participant code, the position number
  (No.\ of the test order) and, after the target, the trial number (for example \code{P01\_03\_knee\_extension\_110deg\_T2\_\dots}); check
  them against the session sheet. Run Export all trials of this participant. \\
13.2 & (add) Never write a name in \lt{}'s Notes field: it is saved in every trial. \\
Test order & Rows 12--15: Hip abduction, hip 0\dg: abduction 0\dg, 20\dg. Hip adduction, hip 0\dg: at abduction 0\dg, 20\dg{} (\lt{}
  shows the abduction angle). Rows 20--21: Belt squat: knee flexion 90\dg, 60\dg. Delete the TBD on the order (\code{protocol.json}
  follows it since 0.20.0) and on adduction 20 (20\dg{} abduction, decided 2026-10-26). \\
Definitions & \emph{Resting force}: the strap force with the participant relaxed in the position (Zero~2, taken before the trial) and, found in the
  recording, the median force from 3.0 to 0.5~s before the push.
  \emph{Push}: where the force was last within 10~\% of its rise above the level before it. \emph{10~\% check}: spread of the peak forces
  of one position over their mean (SOP 10.6). \emph{Tested leg}: the dominant leg, chosen per participant in \lt{}. \emph{limbtrack}: the
  group's camera-based tracking program (\ltversion). \\
Related documents & \lt{} \ltversion{} documentation: README.md, FEATURES.md, SETUP\_AND\_VALIDATION.md, PELVIS\_LANDMARK.md, and this
  document. \\
Revision & Text of 2026-10-29: resting force by Zero~2 instead of the 3~s of rest, reminders instead of stops, foot dots only in
  the ankle tests (Section~\ref{sec:reminders}); the rest of this document still follows 0.20.0. v0.19, 2026-10-26: updated to \lt{} \ltversion: tested leg per participant, resting force 3.0--0.5~s before the push, 10~\%
  check after each trial, test order and adduction at abduction 0/20\dg{} in \code{protocol.json}, belt squat 90/60\dg{} (proposal
  accepted), settings checklist (Table~C1), zero versus resting force for AnyBody, Notes without names, new recordings folder, position
  and trial number in the file names. \\
\bottomrule
\end{longtable}

%=====================================================================================================================================
\part*{Part C. Settings to check at the start of each session}
\addcontentsline{toc}{part}{Part C. Settings to check at the start of each session}

\begin{longtable}{@{}P{4.6cm}P{2.6cm}P{8.8cm}@{}}
\caption*{Table C1: \lt{} settings for the study. \lt{} remembers its settings per computer, so a value can differ from the SOP's.}\\
\toprule
Setting (section) & Value & Note \\
\midrule
\endhead
Tested leg (Participant) & the dominant leg & per participant; Record asks when not chosen \\
Tolerance main joint (Position) & 5.0\dg & default in \code{protocol.json} since 0.20.0 \tbd{10\dg{} after the pilot, if 5\dg{} is not possible} \\
Tolerance other joints (Position) & 5.0\dg & \\
Dot colour (Dots) & Black dots & black felt pads on white tape; \emph{Auto} can follow the tape \\
Smooth the displayed angles (Position) & on & display only \\
Mark the MVIC automatically (Markers) & on & \\
Start when force rises by (Markers) & 25~N & \tbd{confirm; the program's default is 10~N} \\
MVIC length (Markers) & 7.0~s & 1--5--1 \\
Stop by itself after the MVIC end (Markers) & on, 2.0~s & \\
Sound cues (Markers) & on & Test the beeps every session \\
Tick during the hold every (Markers) & 1.0~s & \\
Ruler length (Ruler) & 100~cm & pads 100~cm apart \\
Also save the overlay video (Record) & on & \\
Recordings folder (Record) & the study folder & never the pilot folder \tbd{name} \\
\bottomrule
\end{longtable}

%=====================================================================================================================================
\part*{Part D. Research protocol: sentences to change}
\addcontentsline{toc}{part}{Part D. Research protocol: sentences to change}

\begin{longtable}{@{}P{2.2cm}P{6.0cm}P{7.8cm}@{}}
\toprule
Where & Now & Change to \\
\midrule
\endhead
2.2 & ``National instruments USB-6000''; ``accuracy is $\pm$0.03''; ``A goniometer \dots{} will be used'' & NI USB-6003; the accuracy with
  its unit \tbd{from the data sheet}; ``an electronic angle meter (geo-FENNEL A-Digit 50)''. \\
2.3 & ``the feed will be cropped to exclude the subject's face before recording'' & ``the saved videos are cropped and masked to exclude
  the face; the live picture is used for tracking and is not saved uncropped''. \\
2.3 / MVIC & no trial testing; 1--5--1 & add: one submaximal contraction (about 50~\%) takes up the slack before the trials and is not
  used; the participant rests in the strap for 3~s before each contraction; standardised beeps give the tempo. \\
1.3 & ``two-to-four MVIC measurements in different joint positions'' & ``two or three joint positions per test (Table~2)''. \\
Table 2 caption & ``Dorsal- and plantarflexion are negative'' & ``Hip and knee flexion positive; ankle dorsiflexion positive and
  plantarflexion negative ($-20$\dg{} = 20\dg{} plantarflexion in both ankle tests); hip adduction is tested at 0 and 20\dg{} abduction''. \\
Table 2 & Hip ADD $\theta_2$ = 20\dg & 20\dg{} abduction. \\
Table 3, hip & ``Proximal: lateral midline of the pelvis'' & ``Proximal: from the iliac crest pad (midway between ASIS and PSIS) towards
  the greater trochanter (lateral midline of the pelvis)''. \\
Table 3, pelvic tilt & ``Angle between the ipsilateral ASIS--PSIS line and horizontal'' & remove, or ``observed, not measured''. \\
Analysis & (not stated) & which angle (plateau mean or at peak force) and which force (total zeroed force, see Part~A,
  Section~\ref{sec:anybody}) go to AnyBody; with three trials, which trial counts \tbd{decide}. \\
\bottomrule
\end{longtable}

\begin{thebibliography}{9}
\bibitem{tsaopoulos} D.\,E.\ Tsaopoulos et al., ``Mechanical correction of dynamometer moment for the effects of segment motion during
  isometric knee-extension tests'', \emph{J Appl Physiol} 111(1):68--74, 2011 (SOP reference [45]).
\end{thebibliography}

\end{document}
